\documentclass[12pt,a4paper]{article}

\usepackage[utf8]{inputenc}
\usepackage[T1]{fontenc}
\usepackage[ngerman]{babel}
\usepackage{geometry}
\usepackage{graphicx}
\usepackage{tabularx}
\usepackage{hyperref}
\usepackage{fancyhdr}
\usepackage{enumitem}
\usepackage{xcolor}
\usepackage{colortbl}

\geometry{left=2.5cm, right=2.5cm, top=2.5cm, bottom=2.5cm}

\pagestyle{fancy}
\fancyhf{}
\fancyhead[L]{Lernfeld 7 – Cyber-physische Systeme}
\fancyhead[R]{Arduinoprojekt}
\fancyfoot[C]{\thepage}
\renewcommand{\headrulewidth}{0.4pt}

% Title page info macros
\newcommand{\teilnehmer}{\textcolor{red}{\textbf{[Namen der Teilnehmer]}}}
\newcommand{\datum}{29. Mai 2026}

\begin{document}

% ============================================================
% TITELSEITE
% ============================================================
\begin{titlepage}
\vspace*{2cm}
\begin{center}
    {\LARGE \textbf{Fachinformatiker/in}}\\[0.3cm]
    {\Large Lernfeld 7 – Cyber-physische Systeme ergänzen}\\[0.3cm]
    {\large LS: Durchführung eines Projektes mit Hilfe des Arduinos}\\[2cm]

    {\Huge \textbf{Projekthandout}}\\[0.5cm]
    {\Large Arduino GameBoy – Handheld-Spielekonsole}\\[3cm]

    \begin{tabular}{rl}
        \textbf{Teilnehmer:} & \teilnehmer \\[0.3cm]
        \textbf{Datum:}      & \datum \\[0.3cm]
        \textbf{Klasse:}     & \textcolor{red}{\textbf{[Klasse eintragen]}} \\[0.3cm]
        \textbf{Betreuer:}   & \textcolor{red}{\textbf{[Name der Lehrkraft]}}
    \end{tabular}

    \vfill
    \small Team Informationstechnik (2026)
\end{center}
\end{titlepage}

% ============================================================
% 1. FORMULIERUNG DES THEMAS
% ============================================================
\section{Formulierung des Themas}

\subsection{Thema}
Entwicklung einer Handheld-Spielekonsole auf Basis des Arduino Mega 2560 – nachfolgend
\textit{Arduino GameBoy} genannt. Das Gerät vereint mehrere Spiele in einem kompakten,
tragbaren Gehäuse und wird über ein OLED-Display, einen analogen Joystick sowie zwei
Aktionstasten bedient.

\subsection{Projektidee und Motivation}
Ziel des Projektes ist die Konzeption und Umsetzung eines cyber-physischen Systems, das
die im Unterricht erarbeiteten Grundlagen – Ansteuerung von Displays, Verwendung von
Tastern, Einsatz von Interrupts und Auslesen analoger Sensoren – in einer praktischen
Anwendung bündelt. Die Idee einer Spielekonsole wurde gewählt, weil sie einerseits einen
hohen motivierenden Charakter besitzt und andererseits eine Vielzahl technischer
Herausforderungen bietet: Echtzeitsteuerung, Speicherverwaltung, Zustandsautomaten und
eine modulare Softwarearchitektur.

\subsection{Projektziele}
Die fertige Konsole soll folgende Anforderungen erfüllen:
\begin{itemize}[noitemsep]
    \item \textbf{Homescreen} mit horizontaler Spielauswahl über animierte Karten
    \item \textbf{Mehrere Spiele:} Snake (,,Origin``), Space Invaders (,,Infinity``),
          Dino Jump (,,Phallus``), Vampire Survivors (,,Godly``) und ein
          Raketen-Höhenspiel (,,Eternal``)
    \item \textbf{Pause-Funktion} in jedem Spiel mit Optionen für Weiter, Neustart
          und Zurück zum Menü
    \item \textbf{Highscore-Speicherung} im EEPROM – auch nach Stromverlust erhalten
    \item \textbf{Soundausgabe} über einen aktiven Buzzer für akustisches Feedback
    \item \textbf{Erweiterbare Architektur:} Neue Spiele lassen sich mit geringem
          Aufwand über eine einheitliche Schnittstelle hinzufügen
    \item \textbf{Stabiles 30 FPS Gameplay} trotz begrenzter Rechenleistung des
          Mikrocontrollers
\end{itemize}

% ============================================================
% 2. RESSOURCEN UND ABLAUFPLANUNG
% ============================================================
\section{Ressourcen und Ablaufplanung}

\subsection{Sachmittel und Kostenplanung}

\begin{table}[h]
\centering
\begin{tabularx}{\textwidth}{|l|c|X|}
\hline
\textbf{Komponente}               & \textbf{Anzahl} & \textbf{Kosten (ca.)} \\
\hline
Arduino Mega 2560 (Rev. 3)        & 1 & \textcolor{red}{\textbf{[EUR eintragen]}} \\
2.42\textsuperscript{''} OLED-Display (SSD1309, SPI) & 1 &
    \textcolor{red}{\textbf{[EUR eintragen]}} \\
Joystick-Modul (VRX, VRY, SW)     & 1 & \textcolor{red}{\textbf{[EUR eintragen]}} \\
Taster (Pull-Down)                & 2 & \textcolor{red}{\textbf{[EUR eintragen]}} \\
Aktiver Buzzer                    & 1 & \textcolor{red}{\textbf{[EUR eintragen]}} \\
Widerstände, Jumper-Kabel, Steckbrett & -- & \textcolor{red}{\textbf{[EUR eintragen]}} \\
LEGO-Steine (Gehäuse)             & -- & \textcolor{red}{\textbf{[EUR eintragen]}} \\
\hline
\multicolumn{2}{|r|}{\textbf{Gesamtkosten:}} & \textcolor{red}{\textbf{[EUR eintragen]}} \\
\hline
\end{tabularx}
\caption{Hardware-Komponenten und Kosten}
\end{table}

\subsection{Zeitliche Ablaufplanung}

Der Bearbeitungszeitraum erstreckte sich vom \textcolor{red}{\textbf{[Startdatum]}} bis zum
29.05.2026. Die folgende Tabelle gibt einen Überblick über die wesentlichen Projektphasen.

\vspace{0.3cm}
\begin{tabularx}{\textwidth}{|l|X|l|}
\hline
\textbf{Phase} & \textbf{Tätigkeit} & \textbf{Zeitraum} \\
\hline
1 & Konzeption: Anforderungen definieren, Hardware auswählen,
    Spieleideen skizzieren & \textcolor{red}{\textbf{[KW X]}} \\
2 & Basisaufbau: Display in Betrieb nehmen, Joystick und Taster
    anbinden, Software-Grundgerüst & \textcolor{red}{\textbf{[KW X]}} \\
3 & Spieleentwicklung: Homescreen, Space Invaders, Dino Jump &
    \textcolor{red}{\textbf{[KW X]}} \\
4 & Erweiterung: Vampire Survivors, Snake, Eternal &
    \textcolor{red}{\textbf{[KW X]}} \\
5 & Gehäusebau, Feinabstimmung, Pause-Logic, Highscore-EEPROM &
    \textcolor{red}{\textbf{[KW X]}} \\
6 & Dokumentation, Präsentation, Videoerstellung &
    \textcolor{red}{\textbf{[KW X]}} \\
\hline
\end{tabularx}

% ============================================================
% 3. DURCHFÜHRUNG
% ============================================================
\section{Durchführung}

\subsection{Vorgehensweise und Prozessschritte}

Die Entwicklung folgte einem iterativen Ansatz. Zunächst wurde die grundlegende
Hardware-Umgebung aufgebaut und getestet (Display-Ansteuerung, Eingabe über
Joystick/Taster). Anschließend entstand eine modulare Software-Architektur mit einer
abstrakten \texttt{Game}-Basisklasse, die es erlaubt, Spiele unabhängig voneinander zu
entwickeln und über einen zentralen Zustandsautomaten zu verwalten.

\begin{enumerate}[noitemsep]
    \item \textbf{Hardware-Aufbau:} Verkabelung aller Komponenten gemäß Schaltplan (s.~Anlage).
          Das OLED-Display kommuniziert per Hardware-SPI (MOSI=51, SCK=52), der Joystick
          liefert analoge X/Y-Werte über A0/A1, der Joystick-Taster hängt an Pin~2
          (INPUT\_PULLUP). Zwei externe Taster mit Pull-Down-Widerständen an Pin~30 und~32
          dienen als Aktions- bzw. Pause-Taste. Ein aktiver Buzzer an Pin~40 erzeugt
          akustische Rückmeldungen.
    \item \textbf{Software-Architektur:} Trennung in Hardware-Abstraktion
          (\texttt{hardware.h/.cpp}), Spiele-Basisklasse (\texttt{game\_base.h}),
          Sound-Manager (\texttt{sound.h/.cpp}), Homescreen (\texttt{homescreen.h/.cpp})
          sowie je eine Header/Source-Kombination pro Spiel. Das Hauptprogramm
          (\texttt{GameBoy.ino}) implementiert einen Zwei-Zustands-Automaten
          (Homescreen $\leftrightarrow$ Spiel).
    \item \textbf{Eingabesystem:} Implementierung einer entprellten Button-Klasse mit
          Flankenerkennung (\texttt{consumePress()}). Zusätzlich werden
          Pin-Change-Interrupts (PCINT2 für PORTC, INT4 für Joystick-SW) genutzt, um
          Tastendrücke auch während laufender Spiel-Logik zuverlässig zu erfassen.
    \item \textbf{Display-Rendering:} Verwendung der U8g2-Bibliothek im Full-Framebuffer-Modus.
          Alle Spiele-Grafiken werden mit schnellen Zeichenprimitiven
          (\texttt{drawBox()}, \texttt{drawDisc()}, \texttt{drawXBMP()}) umgesetzt,
          verzögerungsanfällige Operationen wie \texttt{drawLine()} wurden vermieden.
    \item \textbf{Spieleentwicklung:} Jedes Spiel implementiert eine eigene
          Zustandsmaschine (Menü $\rightarrow$ Spielend $\rightarrow$ Pausiert
          $\rightarrow$ Game-Over) und trennt Zeichenmethoden strikt in
          \texttt{drawGame()} (ohne \texttt{sendBuffer()}) und \texttt{drawPlaying()}
          / \texttt{drawPause()} (mit \texttt{sendBuffer()}), um Flackern zu vermeiden.
\end{enumerate}

\subsection{Schaltplan / Visualisierung}

Ein detaillierter Schaltplan mit allen Pin-Verbindungen befindet sich in der Anlage.
Die zentrale Pin-Belegung ist:

\vspace{0.2cm}
\begin{tabularx}{\textwidth}{|l|c|l|X|}
\hline
\textbf{Komponente} & \textbf{Pin} & \textbf{Typ} & \textbf{Hinweis} \\
\hline
OLED CS   & 10 & Digital & SPI Chip-Select \\
OLED DC   &  9 & Digital & SPI Data/Command \\
OLED RST  &  8 & Digital & Display-Reset \\
OLED MOSI & 51 & HW-SPI  & Hardware SPI \\
OLED SCK  & 52 & HW-SPI  & Hardware SPI \\
Joystick X & A0 & Analog  & X-Achse (links/rechts) \\
Joystick Y & A1 & Analog  & Y-Achse (hoch/runter) \\
Joystick SW&  2 & Digital & Taster im Joystick \\
Button 1   & 32 & Digital & Aktion (Schießen/Springen/Bestätigen) \\
Button 2   & 30 & Digital & Pause / Zurück \\
Buzzer     & 40 & Digital & Aktiv (HIGH = ein) \\
\hline
\end{tabularx}

\subsection{Aufgetretene Schwierigkeiten}

Während der Entwicklung traten mehrere Herausforderungen auf:

\begin{itemize}
    \item \textbf{Display zu klein:} Das ursprünglich vorgesehene Display erwies sich als
          zu klein für eine komfortable Spieleanzeige. Es wurde durch ein
          2.42\textsuperscript{''} OLED mit 128×64 Pixeln ersetzt.
    \item \textbf{Arbeitsspeicher unzureichend:} Der Arduino Uno verfügt nur über 2~KB
          SRAM – nicht genug für den Framebuffer (1~KB), mehrere Spielzustände und
          Objekt-Pools. Die Lösung war der Umstieg auf einen Arduino Mega 2560 mit
          8~KB SRAM.
    \item \textbf{Gehäuse-Integration:} Alle Komponenten mit korrekten Pin-Verbindungen
          in ein LEGO-Gehäuse zu integrieren, erforderte mehrere Anläufe, bis alle
          Kabel sauber verlegt und die Taster mechanisch stabil montiert waren.
    \item \textbf{Pause-Funktion:} Das Unterbrechen des laufenden Spielgeschehens und
          das korrekte Wiederaufnehmen erwies sich als komplex. Die Lösung bestand im
          Einsatz von Pin-Change-Interrupts, die einen Pause-Flag setzen, der im
          Hauptloop ausgewertet wird. So wird die Spiel-Logik zuverlässig angehalten,
          während der Pause-Bildschirm weiterhin Eingaben verarbeiten kann.
\end{itemize}

% ============================================================
% 4. PROJEKTERGEBNIS
% ============================================================
\section{Projektergebnis}

\subsection{Soll-Ist-Vergleich}

\begin{tabularx}{\textwidth}{|l|X|X|}
\hline
\textbf{Anforderung} & \textbf{Soll} & \textbf{Ist} \\
\hline
Homescreen        & Spielauswahl mit Vorschau & \textbf{Erfüllt:}
                    Horizontales Karten-Layout mit weicher Scroll-Animation,
                    16×16 Icons, Spielname unter jedem Bild \\
\hline
Spieleanzahl      & Mind. 2 Spiele & \textbf{Übertroffen:}
                    5 Spiele implementiert (Snake, Space Invaders, Dino Jump,
                    Vampire Survivors, Rocket-Launch) \\
\hline
Highscore         & Punktezähler während des Spiels & \textbf{Erfüllt:}
                    Persistente Speicherung im EEPROM pro Spiel, Anzeige im
                    Menü und Game-Over-Bildschirm \\
\hline
Pause / Menü      & Pausieren, Fortsetzen, Neustart & \textbf{Erfüllt:}
                    Debounced Menu-Navigation, 300\,ms Guard gegen
                    versehentliches Doppel-Pausieren \\
\hline
Sound             & Akustisches Feedback & \textbf{Erfüllt:}
                    Nicht-blockierender Buzzer mit \texttt{SoundManager} \\
\hline
Performance       & Flüssiges Gameplay & \textbf{Erfüllt:}
                    Stabile 30\,FPS durch Frame-Timer (33\,ms Intervall),
                    optimierte Zeichenprimitive \\
\hline
Erweiterbarkeit   & Neue Spiele einfach hinzufügbar & \textbf{Erfüllt:}
                    Einheitliche \texttt{Game}-Schnittstelle, neues Spiel in
                    8 Schritten integrierbar \\
\hline
Gehäuse           & Tragbares Gerät & \textbf{Erfüllt:}
                    Kompaktes LEGO-Gehäuse mit allen Komponenten \\
\hline
\end{tabularx}

\subsection{Qualitätskontrolle}

Die Konsole wurde in mehreren Test-Sessions auf Herz und Nieren geprüft:
sämtliche Spiele wurden auf korrektes Verhalten bei Game-Over, Pause und
Menü-Rückkehr getestet. Die Highscore-Persistenz wurde durch Stromtrennung
verifiziert. Das Gehäuse hält auch intensiverem Tastendruck stand, und die
Steckverbindungen bleiben auch bei Bewegung des Gerätes stabil.

\subsection{Abweichungen und Anpassungen}

Die wesentlichen Abweichungen vom ursprünglichen Plan – Umstieg auf ein
größeres Display und den leistungsfähigeren Mega 2560 – haben den Projektumfang
eher erweitert als eingeschränkt. Der zusätzliche Speicher ermöglichte erst die
Umsetzung komplexerer Spiele wie Vampire Survivors und Eternal, die umfangreiche
Objekt-Pools und Zustandsautomaten benötigen.

\subsection{Ausblick}

Das Projekt bietet zahlreiche Anknüpfungspunkte für Weiterentwicklungen:
\begin{itemize}[noitemsep]
    \item \textbf{Weitere Spiele:} Die modulare Architektur erlaubt das einfache
          Hinzufügen von Tetris, Pong oder Jump-\textquotesingle{}n\textquotesingle{}-Run-Varianten.
    \item \textbf{Multiplayer:} Über die serielle Schnittstelle (UART) des Mega
          könnten zwei Geräte für kompetitive Spiele gekoppelt werden.
    \item \textbf{Akku-Betrieb:} Ein Li-Ion-Akku mit Ladeelektronik würde das
          Gerät vollständig mobil machen.
    \item \textbf{Gehäuse-Veredelung:} Ein 3D-gedrucktes, ergonomisch optimiertes
          Gehäuse böte mehr Komfort und eine professionellere Optik.
    \item \textbf{Ton-Verbesserung:} Ein passiver Lautsprecher mit PWM-Ansteuerung
          könnte Melodien statt einfacher Pieptöne erzeugen.
\end{itemize}

% ============================================================
% 5. ANLAGE
% ============================================================
\section{Anlage}

Folgende Materialien sind diesem Handout beigefügt bzw. liegen im Moodle-Abgabeordner
vor:

\begin{enumerate}[noitemsep]
    \item Schaltplan (Fritzing / handgezeichnet)
    \item Vollständiger Programmcode (siehe Moodle-Abgabe)
    \item Screenshots aller Spiele und des Homescreens
    \item Video: Aufbau und Funktion der Konsole (siehe Moodle-Abgabe)
    \item Präsentationsfolien (siehe Moodle-Abgabe)
\end{enumerate}

\end{document}
